\chapter{Static Site Generator (SSG): build a site compiler on top of Relax}
\label{ch:ssg}

\section{What you'll build}
By the end of this chapter, you will have:

\begin{itemize}
  \item A working mental model for how a static site generator works (inputs, pipeline, outputs).
  \item A minimal \emph{SSG pipeline} you can implement yourself: walk files \textrightarrow parse markdown \textrightarrow render HTML \textrightarrow write to disk.
  \item A clear understanding of how Relax.js integrates SSG support \emph{without} changing the browser runtime.
\end{itemize}

\section{What you'll learn}
\begin{itemize}
  \item The core contract of an SSG: deterministic build output from deterministic input.
  \item How Relax's VNode model becomes a server-side HTML renderer (VDOM \textrightarrow string).
  \item Why frontmatter exists, and how Relax parses a "YAML-ish" subset safely.
  \item How clean URLs and sitemap generation work (and why they are just path functions).
  \item How tests act as executable documentation for a build pipeline.
\end{itemize}

\section{Prereqs}
You should have read:

\begin{itemize}
  \item Chapter~4 (VNode model)
  \item Chapter~5 (Mounting)
  \item Chapter~8 (Patching / diffing) --- not because the SSG diffs DOM, but because it explains identity and escaping.
  \item Chapter~23 (Release-quality discipline) helps for the "tests as specs" mindset.
\end{itemize}

\section{Concept: what an SSG really is}
A static site generator is a compiler.
It takes source files (often Markdown) and produces a directory of HTML files.

\subsection{A mental model you can run in your head}
When people say "SSG", they sometimes mean 10 different things (framework, templating system, router, asset pipeline).
In this chapter we are deliberately strict:

\begin{quote}
An SSG is a \textbf{deterministic build pipeline} that turns source content into a folder of deployable files.
\end{quote}

That gives you a clear contract.
If you keep this contract, the tool stays reliable and easy to test.
If you violate it (randomness, time-dependent output, hidden I/O during rendering), you get "works on my machine" builds.

\subsection{The compiler analogy (and why it matters)}
If you've ever used a language compiler, this will feel familiar:

\begin{itemize}
  \item Source code \textrightarrow parse 
  \item AST \textrightarrow transform
  \item IR \textrightarrow codegen
  \item Output binaries
\end{itemize}

An SSG is the same shape:

\begin{itemize}
  \item Markdown files \textrightarrow parse frontmatter + markdown
  \item A "page model" (content + metadata) \textrightarrow wrap in layout
  \item VNodes \textrightarrow HTML string
  \item Output files (HTML, CSS, assets, sitemap)
\end{itemize}

This analogy is useful because compilers are traditionally tested with "golden" outputs and end-to-end fixtures.
That's exactly how we test the Relax SSG.

The useful mental model is:

\begin{quote}
\textbf{SSG contract:} \emph{(content + templates + config) \textrightarrow deterministic files}\\
Same inputs should produce the same output bytes.
\end{quote}

That contract sounds obvious, but it forces good engineering:

\begin{itemize}
  \item you must have predictable path mapping (``where does this page go?'')
  \item you must have predictable rendering (``how does Markdown become HTML?'')
  \item you must have predictable integration points (``how can users customize layout?'')
\end{itemize}

In Relax.js, the SSG is intentionally minimal.
It is not a framework.
It is a sequencing of small, testable parts:

\begin{enumerate}
  \item \texttt{walkFiles()} enumerates input.
  \item \texttt{markdownToPage()} parses content + frontmatter.
  \item \texttt{renderPageToString()} turns VNodes into a full HTML document.
  \item \texttt{toOutputPath()} decides where each page is written.
  \item \texttt{copyDir()} copies public assets.
  \item \texttt{renderSitemapXml()} emits \texttt{sitemap.xml} if you provide a base URL.
\end{enumerate}

\section{Build it from scratch: the smallest pipeline}
Let's build the pipeline conceptually in four steps.
I'll show simplified code first, then point you at the real implementation.

\subsection{Pipeline diagram (map the code before you read it)}
Before we touch details, here is the whole pipeline in one diagram.
Keep this as your "table of contents" for the SSG implementation.

\begin{verbatim}
    (inputDir)
        |
        v
  [ walkFiles() ]  -----------.
        |                     |
        v                     |
      *.md files only               |
        |                     |
        v                     |
     [ markdownToPage() ]           |
       |             |              |
       |             '-> frontmatter (title, layout, etc)
       v
      VNode (main{ innerHTML: ... })
        |
        v
     [ wrapWithLayout() ]  (optional)
        |
        v
     [ renderPageToString() ]
        |
        v
    HTML string  -----> writeFile(toOutputPath())

   plus:
   - copyDir(public -> out)
   - write theme/default.css
   - renderSitemapXml(baseUrl, pages)
\end{verbatim}

Every box corresponds to a real function in \texttt{src/ssg/}.
If you ever want to add a feature, start by deciding \emph{which box it belongs in}.

\subsection{Step 1: walk the input tree}
\textbf{Goal:} take an input directory and return a list of files.

\textbf{Invariant:} the walker must visit every file exactly once.

The real implementation is in \texttt{src/ssg/site.ts} as \texttt{walkFiles()}.
Here's the core idea (simplified):

\begin{lstlisting}[style=relax,language=JavaScript,caption={Walking the input directory (simplified).}]
async function walkFiles(dir) {
  const entries = await readdir(dir)
  const out = []

  for (const name of entries) {
    const abs = join(dir, name)
    const st = await stat(abs)
    if (st.isDirectory()) out.push(...(await walkFiles(abs)))
    else out.push(abs)
  }

  return out
}
\end{lstlisting}

\textbf{Why this exists:}
SSG is "pure" only if input enumeration is deterministic.
If you miss files or include temporary build artifacts, rebuilds become unreliable.

\textbf{What breaks if you change it:}
If the walker accidentally follows output directories, you get exponential growth (page outputs treated as inputs).

\subsection{Determinism note: filesystem order}
Different filesystems can return directory entries in different orders.
If you want perfectly stable output ordering (for stable sitemaps, predictable logs, reproducible builds), a common trick is to sort the directory entries.
Relax's current implementation is intentionally minimal.
If you extend it, consider adding sorting in \texttt{walkFiles()} and covering it with a unit test.

\subsection{Step 2: Markdown \textrightarrow VNodes (but safely)}
Relax's SSG uses Markdown as a content authoring format.
The core pipeline is:

\begin{quote}
Markdown string \textrightarrow HTML string \textrightarrow VNode with \texttt{innerHTML}
\end{quote}

The key file is \texttt{src/ssg/markdown.ts}.
The core function is \texttt{markdownToVNode(markdown)}.

\begin{lstlisting}[style=relax,language=JavaScript,caption={Markdown to VNode uses \texttt{innerHTML} (from \texttt{src/ssg/markdown.ts}).}]
export async function markdownToVNode(markdown, options = {}) {
  const wrapperTag = options.wrapperTag ?? 'div'
  const { marked } = await import('marked')
  const html = marked.parse(markdown)
  return h(wrapperTag, { innerHTML: String(html) }, [])
}
\end{lstlisting}

\subsubsection{Code walk-through: what \texttt{innerHTML} means in a VNode}
In the browser runtime, setting \texttt{innerHTML} is a classic foot-gun.
In an SSG, it's a pragmatic escape hatch.

The important detail is that Relax's string renderer understands this prop.
In \texttt{src/ssg/render-to-string.ts}, element rendering does:

\begin{itemize}
  \item if \texttt{props.innerHTML} is a string: emit it verbatim as children
  \item otherwise: render the VNode children recursively (and escape text nodes)
\end{itemize}

That separates the two worlds clearly: \emph{"safe by default"} vs \emph{"explicit raw"}.

\textbf{Why this design:}
I want the SSG to reuse Relax's VNode model as a "template language" without re-implementing an HTML AST.
The VNode is just a container, and \texttt{innerHTML} carries the rendered markup.

\textbf{Security note:}
For an SSG, Markdown is trusted input (your repo).
If you generate pages from untrusted user input, treat Markdown-to-HTML as a sanitizer problem.
Relax does not attempt to sanitize.

\subsubsection{Practical tip: keep markdown rendering isolated}
Notice how \texttt{marked} is imported lazily:

\begin{itemize}
  \item it keeps the core runtime smaller
  \item it lets consumers of Relax avoid a markdown dependency unless they actually use the SSG
\end{itemize}

This pattern is common in CLI tools: heavy dependencies live behind feature gates.

\subsection{Step 3: Frontmatter is just structured metadata}
Frontmatter exists because content needs metadata:

\begin{itemize}
  \item a title for \texttt{<title>} and headers
  \item a description
  \item a layout toggle
  \item flags like \texttt{draft: true}
\end{itemize}

Relax parses a tiny subset: a flat list of \texttt{key: value} pairs.
This is intentional.

\begin{quote}
\textbf{Design choice:} prefer a small, predictable parser over a "full YAML" dependency.
\end{quote}

The parser lives in \texttt{splitFrontmatter()} in \texttt{src/ssg/markdown.ts}.
Here is the essence:

\begin{lstlisting}[style=relax,language=JavaScript,caption={YAML-ish frontmatter parsing (simplified).}]
function splitFrontmatter(input) {
  const trimmed = input.replace(/^\ufeff/, '')
  if (!trimmed.startsWith('---')) return { content: input, data: {} }

  const end = trimmed.indexOf('\n---', 3)
  if (end === -1) return { content: input, data: {} }

  const fmBlock = trimmed.slice(3, end).trim()
  const rest = trimmed.slice(end + 4)

  const data = {}
  for (const line of fmBlock.split(/\r?\n/)) {
    const m = /^([A-Za-z0-9_\-]+)\s*:\s*(.*)$/.exec(line)
    if (!m) continue
    data[m[1]] = parseScalar(m[2])
  }

  return { content: rest.replace(/^\s*\r?\n/, ''), data }
}
\end{lstlisting}

\subsubsection{Reading the real implementation: what it deliberately does \emph{not} do}
The Relax frontmatter parser is intentionally small.
When you read \texttt{splitFrontmatter()} in \texttt{src/ssg/markdown.ts}, you should notice:

\begin{itemize}
  \item It strips a UTF-8 BOM (byte-order mark). This is a real-world paper-cut.
  \item It only parses lines that look like \texttt{key: value}.
  \item It parses scalars: \texttt{true/false}, numbers, and strings (including quoted strings).
  \item It ignores nested YAML, lists, multiline blocks, anchors, and other YAML features.
\end{itemize}

This is a feature, not a limitation.
It keeps the parser predictable and reduces surprises in a tutorial.

\subsubsection{Diagram: how a markdown file becomes a \texttt{Page}}
\begin{verbatim}
index.md
  |
  |  splitFrontmatter()
  v
{ data: { title: "Home" }, content: "# Hello" }
  |
  |  markdownToVNode()
  v
{ vdom: h('main', { innerHTML: "<h1>Hello</h1>" }), data }
\end{verbatim}

\textbf{Why this exists:}
Frontmatter is how content communicates with the renderer without having to embed logic in Markdown.

\textbf{What breaks if you remove it:}
You lose page titles, descriptions, and any ability to make layout decisions at build time.

\subsection{Step 4: VNodes \textrightarrow HTML string (SSR-friendly rendering)}
To generate HTML files, we need a server-side renderer.
This is not DOM mount.
This is pure string rendering.

Relax implements this in \texttt{src/ssg/render-to-string.ts}.

The important contract:

\begin{quote}
\texttt{renderToString(vdom)} must escape by default, but allow controlled raw HTML for Markdown.
\end{quote}

The key design is \texttt{innerHTML}:

\begin{itemize}
  \item normal text children are escaped to prevent accidental tag injection
  \item \texttt{innerHTML} bypasses escaping because Markdown already produced HTML
\end{itemize}

You can see this in the element renderer:

\begin{lstlisting}[style=relax,language=JavaScript,caption={\texttt{innerHTML} rendering in \texttt{renderToString()}.}]
function renderElementToString(vdom) {
  const tag = vdom.tag
  const attrs = renderAttrs(vdom.props ?? {})

  const innerHTML = (vdom.props ?? {}).innerHTML
  const children =
    typeof innerHTML === 'string'
      ? innerHTML
      : (vdom.children ?? []).map(renderToString).join('')

  return `<${tag}${attrs}>${children}</${tag}>`
}
\end{lstlisting}

\subsubsection{Deep dive: escaping rules (the two invariants)}
There are two invariants that make this safe enough for "content from your repo":

\begin{enumerate}
  \item \textbf{Text nodes are always escaped.} If you render a text child like \texttt{"<b>"}, you will get \texttt{\&lt;b\&gt;}.
  \item \textbf{Raw HTML must be explicit.} Only \texttt{innerHTML} bypasses escaping.
\end{enumerate}

You can find these invariants in tests in \texttt{src/\_\_tests\_\_/ssg.test.js}.
In build tools, tests are not just correctness checks --- they are documentation of your security posture.

\subsubsection{What about components?}
Relax's renderer can also render component VNodes.
This is useful if you later want to build small "islands" of reusable layout logic (headers, footers, callouts) as components.
The renderer instantiates and calls \texttt{render()} without mounting.

This is a powerful idea:

\begin{quote}
You can reuse your UI component model as a static templating system.
\end{quote}

It also has a limitation: because nothing is mounted, lifecycle hooks and browser APIs do not run.
That's exactly what we want in an SSG.

\textbf{Why this design:}
In an SSG, performance comes from having a boring, predictable renderer.
This avoids having to create a DOM (jsdom) just to produce HTML.

\section{Integration: the build orchestrator (CLI / API)}
Now let's connect those parts into a build function.
The public entry point today is \texttt{buildStaticSite()} in \texttt{src/ssg/cli.ts}.

\subsection{Build orchestrator: read it like a story}
Open \texttt{src/ssg/cli.ts} and skim \texttt{buildStaticSite()} once without reading every line.
Try to identify these phases:

\begin{enumerate}
  \item Resolve paths (input/output)
  \item Load and merge config
  \item Enumerate markdown files
  \item For each markdown file: parse \textrightarrow wrap \textrightarrow render \textrightarrow write
  \item Write shared outputs (theme, public assets, sitemap)
\end{enumerate}

Then come back and read it slowly.
This skill ("read code by finding the phases") is how you get productive with build tools quickly.

The pipeline is:

\begin{enumerate}
  \item walk input files
  \item filter \texttt{.md}
  \item render each Markdown file to an HTML page
  \item choose its output path (clean URL mapping)
  \item write the file
  \item add URL to sitemap list
  \item copy theme + public assets
  \item write \texttt{sitemap.xml} if configured
\end{enumerate}

\subsection{Clean URLs are just path rules}
Clean URLs are the difference between:

\begin{itemize}
  \item \texttt{about.html} (file URL)
  \item \texttt{about/index.html} (directory URL, easier for static hosts)
\end{itemize}

Relax's mapping is implemented in \texttt{toOutputPath()}:

\begin{lstlisting}[style=relax,language=JavaScript,caption={Clean URL mapping (from \texttt{src/ssg/site.ts}).}]
// index.md -> index.html always
// about.md -> about/index.html when cleanUrls === 'always'
export function toOutputPath(relMdPath, cleanUrls) {
  if (relMdPath === 'index.md') return 'index.html'
  return cleanUrls === 'always'
    ? stripExt(relMdPath) + '/index.html'
    : relMdPath.replace(/\.md$/i, '.html')
}
\end{lstlisting}

\textbf{Why this design:}
The mapping is pure and unit-testable.
It doesn't touch the filesystem.
That keeps mistakes from becoming production bugs.

\subsubsection{Tiny correctness checklist for URL mapping}
When you work on clean URLs, there are 4 easy-to-miss cases:

\begin{itemize}
  \item \texttt{index.md} at the root
  \item nested \texttt{docs/index.md}
  \item Windows path separators (\texttt{\\})
  \item links in sitemap must use URL slashes, not OS separators
\end{itemize}

Relax's \texttt{toOutputPath()} and \texttt{toUrlFromOutPath()} normalize separators to forward slashes.
This is a classic "cross platform" trap, and it is worth testing.

\subsection{Sitemap generation is just "list of pages" \textrightarrow XML}
Sitemaps are another example where a small pure function is ideal:

\begin{quote}
entries \textrightarrow XML string
\end{quote}

Relax does this in \texttt{renderSitemapXml(baseUrl, entries)}.

\section{Config: making a build tool usable}
Build tools need configuration.
Relax supports a project-level config file:

\begin{quote}
\texttt{relax.ssg.config.js}
\end{quote}

The loader lives in \texttt{src/ssg/config.ts}:

\begin{lstlisting}[style=relax,language=JavaScript,caption={Loading \texttt{relax.ssg.config.js} via dynamic import (from \texttt{src/ssg/config.ts}).}]
export async function loadSsgConfig(cwd, configPath) {
  const abs = resolve(cwd, configPath ?? 'relax.ssg.config.js')
  try {
    await readFile(abs, 'utf8')
  } catch {
    return { path: null, config: {} }
  }

  const mod = await import(pathToFileURL(abs).href)
  const config = mod?.default ?? mod ?? {}
  return { path: abs, config }
}
\end{lstlisting}

\subsection{Deep dive: config merging and override rules}
Relax merges three sources of truth:

\begin{enumerate}
  \item built-in defaults (\texttt{DEFAULT\_CONFIG})
  \item project config file (\texttt{relax.ssg.config.js})
  \item CLI / environment overrides (the options passed to \texttt{buildStaticSite()})
\end{enumerate}

The merge rule is simple:

\begin{quote}
	extbf{overrides win}, but we keep base nav if you don't override it.
\end{quote}

That last part matters.
It means a config file can declare a nav once, and callers can still override other fields for CI or preview builds.

This behavior is locked in by \texttt{src/\_\_tests\_\_/ssg-config.test.ts}.

\subsection{Note on build environments (Node vs test runners)}
Build tools often run in strange environments:

\begin{itemize}
  \item Node directly (CLI)
  \item bundled execution (Rollup)
  \item test runners that transform modules (Vitest/Vite)
\end{itemize}

The practical lesson: keep configuration loading boring and test it.
If you ever see a build pass locally but fail in CI, configuration loading and path resolution are top suspects.

\textbf{Why this design:}
A JavaScript config file keeps the SSG friction low:

\begin{itemize}
  \item it can compute values
  \item it can reuse constants
  \item it doesn't require us to ship a YAML parser
\end{itemize}

\subsection{Config-driven navigation (\texttt{nav})}
The default layout wrapper includes a minimal navigation bar.
Hard-coding links is fine for a demo, but a real SSG needs a declarative nav model.

In Relax's SSG config, \texttt{nav} is an array of:

\begin{quote}
	exttt{\{ label: string, href: string \}}
\end{quote}

Example \texttt{relax.ssg.config.js}:

\begin{lstlisting}[style=relax,language=JavaScript,caption={Config nav entries drive the rendered header nav.}]
export default {
  siteName: 'My Site',
  nav: [
    { label: 'Home', href: '/' },
    { label: 'Docs', href: '/docs/' },
  ],
}
\end{lstlisting}

The CLI's default wrapper (\texttt{src/ssg/cli.ts}) maps \texttt{nav} to links:

\begin{lstlisting}[style=relax,language=JavaScript,caption={Default layout nav rendering (simplified).}]
const navItems = ctx.nav && ctx.nav.length
  ? ctx.nav
  : [{ label: 'Home', href: '/' }]

h('nav', { class: 'site-nav' }, [
  ...navItems.map((item) => h('a', { href: item.href }, [item.label]))
])
\end{lstlisting}

	extbf{Invariant:} navigation rendering must be pure.
No filesystem reads.
No global state.
It should only depend on the merged config and frontmatter.

	extbf{What breaks if you violate it:}
the same content build can produce different output HTML, which makes caching and CI diffs miserable.

\subsubsection{Diagram: config \textrightarrow layout context \textrightarrow HTML}
\begin{verbatim}
relax.ssg.config.js
  nav: [{label:'Home',href:'/'},{label:'Docs',href:'/docs/'}]
          |
          v
 mergeConfig()
          |
          v
 wrapWithLayout(ctx={ siteName, nav })
          |
          v
 <nav>
   <a href="/">Home</a>
   <a href="/docs/">Docs</a>
 </nav>
\end{verbatim}

\section{Tests as specs (SSG correctness)}
The SSG is small enough that tests are its documentation.
There are two layers:

\subsection{Unit tests: pure functions + invariants}
Look at \texttt{src/\_\_tests\_\_/ssg.test.js}.
It proves key invariants:

\begin{itemize}
  \item frontmatter parsing handles types (string/number/bool)
  \item BOM handling (byte-order mark) doesn't break parsing
  \item escaping behavior is correct (text nodes escaped, \texttt{innerHTML} not escaped)
  \item clean URL mapping produces the right output paths and URLs
  \item sitemap XML contains expected URLs
\end{itemize}

A representative assertion (escaping vs raw HTML):

\begin{lstlisting}[style=relax,language=JavaScript,caption={Escaping is the default; \texttt{innerHTML} is an explicit escape hatch.}]
const vdom = h('p', {}, ['<script>'])
expect(renderToString(vdom)).toBe('<p>&lt;script&gt;</p>')

const raw = h('div', { innerHTML: '<span>ok</span>' })
expect(renderToString(raw)).toBe('<div><span>ok</span></div>')
\end{lstlisting}

\subsection{Integration test: "build a site" end to end}
The integration spec is \texttt{src/\_\_tests\_\_/ssg-build.integration.test.ts}.
It creates a temp directory, writes markdown + a \texttt{public/} asset, runs \texttt{buildStaticSite()}, then asserts:

\begin{itemize}
  \item HTML files exist at the expected clean URLs
  \item theme CSS is written
  \item public assets are copied
  \item sitemap contains baseUrl + pages
\end{itemize}

This is the kind of test you want for a build tool.
It prevents accidental regressions where the build "succeeds" but produces missing files.

\subsubsection{How to design good integration fixtures}
The integration test is intentionally tiny.
It creates:

\begin{itemize}
  \item 2 markdown pages
  \item 1 public asset
  \item 1 config file with a nav entry
\end{itemize}

This is enough to cover multiple subsystems without making the test slow.
For build tools, a good heuristic is:

\begin{quote}
	extbf{Use the smallest fixture that exercises multiple outputs.}
\end{quote}

\subsection{Config merge test: predictable overrides}
Config is also unit-tested in \texttt{src/\_\_tests\_\_/ssg-config.test.ts}.
It proves that:

\begin{itemize}
  \item overrides win when provided
  \item base \texttt{nav} is preserved if you don't override it
\end{itemize}

\section{Why this design? (tradeoffs and alternatives)}
This SSG intentionally avoids becoming a second framework.
Here are the key tradeoffs I made:

\subsection{Tradeoff 1: Markdown as HTML-in-a-box}
\textbf{Choice:} Markdown is converted to an HTML string and embedded via \texttt{innerHTML}.

\textbf{Alternative:} Parse markdown into an AST and map it into VNodes directly.

\textbf{Why I didn't:}
AST-to-VDOM is "purer" but it forces you to maintain a mapping layer and deal with edge cases (tables, code fences, custom HTML).
The goal here is to show the pipeline, not build a full markdown renderer.

\subsection{Tradeoff 2: Tiny frontmatter parser}
\textbf{Choice:} support only flat \texttt{key: value} pairs.

\textbf{Alternative:} full YAML frontmatter.

\textbf{Why I didn't:}
Full YAML parsing adds a dependency and a larger attack surface.
For a tutorial-first SSG, most sites only need a handful of scalar fields.

\subsection{Tradeoff 3: One default layout}
\textbf{Choice:} provide a built-in wrapper layout as a teaching example.

\textbf{Alternative:} user-provided templates on disk.

\textbf{Why I didn't (yet):}
File-based templating is a second system (loader + watch + partials + caching).
Keeping a single built-in layout lets the book focus on the build pipeline.

\section{Common pitfalls}
\begin{itemize}
  \item \textbf{Forgetting escaping:} never treat raw strings as HTML unless you explicitly opt into \texttt{innerHTML}.
  \item \textbf{Mixing input and output trees:} don't write build output into the input directory.
  \item \textbf{Path normalization:} always normalize separators when converting filesystem paths into URLs.
  \item \textbf{Non-deterministic paths:} keep path mapping pure and unit-tested (\texttt{toOutputPath}).
  \item \textbf{Missing baseUrl:} sitemap generation needs an absolute base URL.
  \item \textbf{Assuming a browser:} SSG runs in Node, so avoid DOM APIs.
\end{itemize}

\section{Extension ideas (when you're ready)}
If you want to turn this minimal SSG into a more practical tool, here are extensions that fit neatly into the existing boxes:

\begin{itemize}
  \item \textbf{Draft pages:} \texttt{draft: true} in frontmatter skips writing and excludes from sitemap.
  \item \textbf{Collections:} treat \texttt{posts/*.md} as a list, generate index pages.
  \item \textbf{Syntax highlighting:} run a highlighter on Markdown code fences before \texttt{innerHTML}.
  \item \textbf{RSS feed:} similar to sitemap, just a different serializer.
  \item \textbf{Pretty HTML:} add an option to format the output (tradeoff: less deterministic whitespace).
  \item \textbf{Watch mode:} incremental rebuild on file changes (be careful with correctness).
\end{itemize}

\section{Mini-exercise}
Add a new frontmatter key called \texttt{draft} (boolean).
Then modify the build pipeline so drafts are not written to disk or included in \texttt{sitemap.xml}.

\textbf{How to approach it (test-first):}

\begin{itemize}
  \item Add a unit test to \texttt{src/\_\_tests\_\_/ssg-build.integration.test.ts}:
    create a \texttt{draft.md} with \texttt{draft: true} and assert it doesn't exist in output.
  \item Update \texttt{buildStaticSite()} to skip pages where \texttt{page.data.draft === true}.
\end{itemize}

This exercise is small but practical: it forces you to treat metadata as part of the build contract.

\subsection{Bonus exercise: add a page-level layout toggle}
Today, frontmatter can set \texttt{layout: false} to disable wrapping.
Add a new option \texttt{layout: "default" | false | "minimal"} and implement an extra minimal wrapper.

	extbf{Hint:}
keep wrappers as pure VNode functions and unit-test them by asserting for key substrings in the output HTML.
